\documentclass[11pt]{article}
\usepackage[a4paper,margin=24mm]{geometry}
\usepackage{amsmath,amssymb,amsthm,mathtools,hyperref,listings}
\newtheorem{proposition}{Proposition}
\newtheorem{corollary}[proposition]{Corollary}
\newcommand{\Ric}{\operatorname{Ric}}
\newcommand{\tr}{\operatorname{tr}}
\newcommand{\Sym}{\operatorname{Sym}}
\newcommand{\vol}{\,d\mu_g}
\lstset{basicstyle=\ttfamily\scriptsize,breaklines=true,columns=fullflexible,keepspaces=true}
\title{Negative Goldberg: torsion-zero indices and the full Weyl equation}
\author{Research continuation}
\date{15 October 2026}
\begin{document}
\maketitle
\begin{abstract}
The target is the closed four-dimensional negative Goldberg conjecture. From the torsion-zero equations we derive positivity of transverse local indices and exclude compact clean two-dimensional zero components when the Einstein scalar curvature is nonzero. We also give an exact rational almost-K\"ahler metric germ satisfying $\Ric+g=O(|x|^6)$ with an isolated transverse torsion zero, and derive the complete projected Einstein Weyl equation without dividing by torsion. These results do not prove the conjecture or produce a closed counterexample. No claim of historical priority is made for consequences of classical identities.
\end{abstract}

\section{The fixed problem and conventions}
Let $(M^4,g,\omega,J)$ be a connected almost-K\"ahler manifold, oriented by $\omega^2$, with $|\omega|^2=2$ and $\Ric=(s/4)g$, where $s<0$ is constant. The conjecture asks whether compactness without boundary forces $\nabla\omega=0$. We use the operator norm on $W^+:\Lambda^+\to\Lambda^+$ and the nonnegative rough Laplacian $\nabla^*\nabla$.

Write $u=|\nabla\omega|^2$, $e=\omega/\sqrt2$, $E=e^\perp\subset\Lambda^+$, and
\begin{equation}\label{block}
W^+=\begin{pmatrix}a_0&\zeta^T\\\zeta&-a_0I_E/2+B\end{pmatrix},\qquad
 a_0=\frac{s}{6}+\frac{u}{4},\quad B\in\Sym^2_0E,\quad \eta=\sqrt2\zeta.
\end{equation}
Thus $\eta=(W^+\omega)^\perp$. The component $B$ is the Hermitian-anti-invariant Weyl component usually denoted $W_3^+$ or $W''$; it is different from $\eta$.

Armstrong's characteristic-class obstruction gives
\begin{equation}\label{chern}
c_2(\Lambda^{0,1}\otimes\Lambda^{0,2})[M]
=2c_1^2+c_2=5\chi+6\tau.
\end{equation}
For a closed Einstein four-manifold with $s\ne0$, Chern--Weil gives $\chi>0$ and $2\chi+3\tau>0$, hence $5\chi+6\tau=\chi+2(2\chi+3\tau)>0$. The torsion therefore has a zero. The classical one-zero propagation theorem additionally assumes $B\equiv0$; see \cite[Theorem 3 and the preceding argument]{AD}. Neither the existence of a zero nor Einstein analyticity removes that hypothesis.

\section{The actual first-derivative constraint at a zero}
Choose an adapted orthonormal coframe at a torsion zero $o$ so that
\[
\omega=e^{12}+e^{34},\quad \phi=e^{13}-e^{24},\quad \psi=e^{14}+e^{23}.
\]
Write locally
\[
\nabla\omega=a\otimes\phi-Ja\otimes\psi.
\]
The real one-form $a$ vanishes precisely when the Nijenhuis tensor vanishes, and $u=4|a|^2$. At $o$, its derivative is independent of the choice of connection and of derivatives of the adapted frame. Identify both the tangent and coefficient spaces with $\mathbb C^2$ using $(x_1+ix_2,x_3+ix_4)$ and $(a_1+ia_2,a_3+ia_4)$.

\begin{proposition}\label{index}
At a torsion zero of an Einstein almost-K\"ahler four-manifold, the real derivative $A=Da(o)$ has the form
\begin{equation}\label{linear}
 Az=c z+Q\bar z,\qquad c\in\mathbb C,\quad Q\in M_2(\mathbb C),
\end{equation}
and satisfies
\begin{equation}\label{constraint}
\tr(Q\bar Q)=-8|c|^2.
\end{equation}
Consequently
\begin{equation}\label{det}
\det_{\mathbb R}A=9|c|^4+|\det_{\mathbb C}Q|^2\ge0.
\end{equation}
In particular, every transverse torsion zero has local index $+1$.
\end{proposition}
\begin{proof}
The Einstein condition makes the curvature of the Levi-Civita connection on $\Lambda^+$ self-dual. At $a=0$, its two transverse curvature equations state that $da$ and $d(Ja)$ are self-dual. Splitting an arbitrary real map into $Cz+Q\bar z$, these six real equations give $C=cI_2$. This proves \eqref{linear}. The same curvature equations, with the above two-form normalization, give
\begin{equation}\label{norms}
|A|^2=4|c|^2+2\|Q\|_{\mathrm{HS}}^2,\quad
|\eta|^2=8|c|^2,\quad |B|^2=2|Q_{12}-Q_{21}|^2.
\end{equation}
All three formulas are pointwise at $o$.

To obtain the next compatibility condition, evaluate the classical Einstein Sekigawa identity at $o$ \cite[Proposition 2]{AD}. Terms containing $u$, $u^2$, or $\phi_{\mathrm{Sek}}$ vanish. Since $\nabla\omega(o)=0$,
\[
\Delta u=-2|\nabla^2\omega|^2,\qquad
\delta\langle\rho^*,\nabla_{\cdot}\omega\rangle
=\langle\rho^*,\nabla^*\nabla\omega\rangle=2|\eta|^2.
\]
Here $\nabla^*\nabla\omega=2\eta$ at $o$, by the harmonic two-form Weitzenb\"ock formula. The surviving terms of Sekigawa's identity therefore give
\[
16|\eta|^2+2|\nabla^2\omega|^2=8|B|^2+4|\eta|^2.
\]
As $|\nabla^2\omega|^2=4|A|^2$, this is
\begin{equation}\label{pointnorm}
|A|^2+\frac32|\eta|^2-|B|^2=0.
\end{equation}
Substitution of \eqref{norms} gives
\[
16|c|^2+2\bigl(\|Q\|_{\mathrm{HS}}^2-|Q_{12}-Q_{21}|^2\bigr)=0.
\]
The expression in parentheses is $\tr(Q\bar Q)$, proving \eqref{constraint}.
Finally the real determinant of \eqref{linear} equals
\[
\det\begin{pmatrix}cI&Q\\\bar Q&\bar cI\end{pmatrix}
=|c|^4-|c|^2\tr(Q\bar Q)+|\det Q|^2.
\]
This polynomial identity remains valid at $c=0$. Equation \eqref{det} follows. To fix the complex bundle convention explicitly, put $Z_1=e_1-ie_2$, $Z_2=e_3-ie_4$. Direct evaluation of the Nijenhuis tensor, without the optional factor $1/4$ in its definition, gives
\[
N_J(\bar Z_1,\bar Z_2)
=4i\bigl((a_1+ia_2)Z_1+(a_3+ia_4)Z_2\bigr).
\]
Thus the coefficient identification preserves the real orientation of the complex rank-two bundle $T^{1,0}M\otimes\bar K$, where $K=\Lambda^{2,0}T^*M$. This is the torsion bundle in \eqref{chern}.
\end{proof}

\begin{corollary}
If the torsion of a closed Einstein almost-K\"ahler four-manifold is transverse to the zero section, then
\[
\#\{N_J=0\}=5\chi+6\tau.
\]
\end{corollary}
This is a conditional count, not a transversality theorem. Einstein structures cannot be perturbed arbitrarily while preserving their equations. In particular the positive characteristic number does not contradict the allowed local indices.

\section{A global exclusion for clean zero curves}
At a singular zero, \eqref{det} gives $c=0$ and $\det Q=0$. Hence the only possible real ranks of $A$ are $0,2,4$. In rank two, write $Q=v w^T$. Then
\[
0=\tr(Q\bar Q)=|w^T\bar v|^2,
\qquad Q\bar Q=0.
\]
Consequently $A^2=0$, and $\operatorname{im}A=\ker A$ is a complex line. These statements are invariant under changes of the adapted frame, including its phase.

\begin{proposition}\label{clean}
An Einstein almost-K\"ahler four-manifold with $s\ne0$ has no compact embedded surface $C$ contained in $\{N_J=0\}$ such that
\[
\ker (DN_J)_x=T_xC\qquad\text{for every }x\in C.
\]
Here $C$ has no boundary. In particular, there is no compact clean two-dimensional component of the torsion-zero locus.
\end{proposition}
\begin{proof}
The derivative at a zero is intrinsic. The hypothesis gives rank two, so the preceding algebra implies that $TC$ is $J$-invariant and that the image of $DN_J$ is $T^{1,0}C\otimes\bar K|_C$. Moreover, this derivative is anti-complex-linear in its normal argument. If $\mathcal N_C$ denotes the complex normal line bundle, it therefore induces an isomorphism
\begin{equation}\label{normal}
\overline{\mathcal N_C}\simeq T^{1,0}C\otimes\bar K|_C.
\end{equation}
Put $t=c_1(TC)$, $n=c_1(\mathcal N_C)$ and $\gamma=c_1(TM,J)|_C$. The splitting of $TM|_C$ gives $\gamma=t+n$, whereas \eqref{normal} gives $-n=t+\gamma$. Thus $2\gamma=0$ and
\begin{equation}\label{czerodegree}
\langle c_1(TM,J),[C]\rangle=0.
\end{equation}

For the opposite conclusion, use the canonical Hermitian connection on the anti-canonical line bundle. Its curvature $F$ represents $2\pi c_1(TM,J)$ after multiplication by $i$. On an Einstein almost-K\"ahler four-manifold, Blair's curvature formula is
\[
iF=\eta+\frac{s+s^*}{8}\omega+\frac{s-s^*}{8}\widehat\omega,
\qquad s^*-s=u,
\]
where the last term extends as zero at $u=0$; see \cite[Proposition 2.1]{Minimality}. On $C$, $u=0$, while the $J$-anti-invariant form $\eta$ restricts to zero on the $J$-complex tangent line. Therefore
\[
iF|_C=\frac s4\omega|_C,\qquad
\langle c_1(TM,J),[C]\rangle
=\frac{s}{8\pi}\operatorname{Area}_g(C)\ne0,
\]
contradicting \eqref{czerodegree}.
\end{proof}
The ambient manifold need not be compact for Proposition \ref{clean}; compactness of the zero curve is the essential global hypothesis. The result leaves isolated transverse zeros and non-clean zero sets. In particular, it does not assert that every analytic zero component is clean. Across rank-dropping points, \eqref{normal} need not be an isomorphism, and no holomorphic extension has been proved.

\section{A rational compatible metric realizing the allowed zero to finite order}
For six real functions define
\[
\mathcal H(f_1,\ldots,f_6)=
\begin{pmatrix}
f_1&f_2&f_3&f_4\\f_2&-f_1&f_4&-f_3\\
f_3&f_4&f_5&f_6\\f_4&-f_3&f_6&-f_5
\end{pmatrix}.
\]
Let $\omega_0=dx_1\wedge dx_2+dx_3\wedge dx_4$ and set $G=\exp(H_2+H_4+H_6)$, where $H_d$ is homogeneous of degree $d$, $H_2=\mathcal H(f^{(2)})$ and
\[
f^{(2)}=\left(-2x_2x_3-\tfrac32x_3^2-\tfrac12x_4^2,
-2x_1x_3,-x_1x_3-x_2x_4,-2x_1x_4,0,0\right).
\]
The quartic and sextic polynomials defining $H_4,H_6$ are specified in the executable certificate below. Every $H_d$ is symmetric and anticommutes with the standard symplectic matrix $\Omega_0$. Thus $G\Omega_0G=\Omega_0$, $G$ is positive definite, and $G$ defines an exact real-analytic almost-K\"ahler metric compatible with $\omega_0$.

Direct expansion gives
\begin{equation}\label{finite}
\Ric(G)+G=O(|x|^6),\qquad
 a(x)=(x_1-x_4,-x_2-x_3,x_2/2,x_1/2)+O(|x|^3).
\end{equation}
The Jacobian determinant is $1/4$. Equivalently $c=0$ and
\[
Q=\begin{pmatrix}1&-i\\ i/2&0\end{pmatrix},\quad
\tr(Q\bar Q)=0,\quad |\det Q|^2=1/4.
\]
The actual metric has negative Ricci curvature on a sufficiently small neighborhood and an isolated transverse torsion zero. It is only approximately Einstein. Formula \eqref{finite} establishes compatibility through the displayed Taylor order, not an exact Einstein germ or extension to every order. Independent rational Christoffel computations verified every constant, quadratic and quartic Einstein coefficient; odd coefficients vanish by parity. No formal-integrability or convergence conclusion follows from this finite calculation.

\section{The remaining Weyl component has a coupled equation}
Let $D$ be the projected connection preserving $\mathbb Re\oplus E$. Write
\[
\nabla=D+C,\qquad
C_i=\begin{pmatrix}0&-p_i^T\\p_i&0\end{pmatrix}.
\]
The almost-K\"ahler type identity and harmonicity of $\omega$ give
\begin{equation}\label{ptype}
\sum_i p_ip_i^T=\frac u4 I_E,\qquad D^*p=2\zeta.
\end{equation}
These are smooth across $u=0$.

\begin{proposition}
For an Einstein almost-K\"ahler four-manifold, the block $B$ in \eqref{block} satisfies
\begin{equation}\label{Beq}
D^*DB+\left(\frac{3s}{2}+2u\right)B
=2\sum_i[p_i(D_i\zeta)^T+(D_i\zeta)p_i^T]_0
+2(\zeta\zeta^T)_0.
\end{equation}
On a closed manifold this implies
\begin{equation}\label{Benergy}
\int_M\left[|DB|^2+\left(\frac{3s}{2}+2u\right)|B|^2
+4\sum_i\langle(D_iB)p_i,\zeta\rangle
-10\langle B\zeta,\zeta\rangle\right]\vol=0.
\end{equation}
\end{proposition}
\begin{proof}
Use the Einstein Weyl identity \cite[equation (2), with $f=1$]{LeBrun}
\[
\nabla^*\nabla W^++\frac s2W^+-6(W^+)^2+2|W^+|^2I=0.
\]
At a $D$-normal point the rough Laplacian is
\[
D^*DW^+-2\sum_i[C_i,D_iW^+]
-\sum_i[D_iC_i,W^+]-\sum_i[C_i,[C_i,W^+]].
\]
The trace-free $E$-blocks of the four summands are, respectively,
\[
D^*DB,\quad-2\sum_i[p_i(D_i\zeta)^T+(D_i\zeta)p_i^T]_0,
\quad4(\zeta\zeta^T)_0,\quad\frac u2B.
\]
Also $((W^+)^2)_{E,0}=(\zeta\zeta^T)_0-a_0B$, because $B^2$ is scalar in real rank two. This proves \eqref{Beq}. Pair with $B$, integrate, and use $D^*p=2\zeta$ to obtain \eqref{Benergy}.
\end{proof}
The coefficient $3s/2$ is negative at a torsion zero. The source involving $\zeta$ and its derivative is present in the full Einstein system. Thus \eqref{Beq} is not a homogeneous equation with positive potential. No division by $u$ or singular adapted normalization was used.

\section{What this continuation establishes}
The zero-derivative equations have a global consequence: a compact clean zero curve is impossible for nonzero Einstein scalar curvature. They also rule out a proposed contradiction between positive total Chern number and nonpositive isolated indices: transverse indices are positive. The rational metric shows that the Einstein equations through the tested Taylor order allow such a zero. The full $B$ equation specifies the coupling left out by a standalone Weyl-line-bundle vanishing argument.

The remaining task is still to prove $u\equiv0$ for every closed negative Einstein almost-K\"ahler four-manifold, or construct a closed example with $u\not\equiv0$. Neither a convergent local Einstein germ with the displayed zero nor a new global vanishing estimate is established here. Previously obtained complete noncompact examples and their quotient obstructions are inherited, not recounted as new results.

\appendix
\section{Self-contained exact polynomial certificate}
The following Python 3 program uses only SymPy and standard-library JSON. It checks the compatible matrix constraints, every constant, quadratic and quartic coefficient of $\Ric+g$, and the leading torsion derivative, independently of the linear-system construction. All polynomial coefficients are rational.
\begin{lstlisting}[language=Python]
import json
from pathlib import Path
import sympy as s

D=json.loads(r'''{
  "variables": [
    "x1",
    "x2",
    "x3",
    "x4"
  ],
  "H2": [
    "-2*x2*x3 - 3*x3**2/2 - x4**2/2",
    "-2*x1*x3",
    "-x1*x3 - x2*x4",
    "-2*x1*x4",
    "0",
    "0"
  ],
  "H4": [
    "-x1**3*x4/6 + 21*x1**2*x2**2/4 - 57*x1**2*x2*x3/16 + 179*x1**2*x3**2/32 - 267*x1**2*x4**2/32 - 7*x1*x2**2*x4/2 + 9*x1*x2*x3*x4 - 83*x1*x3**2*x4/4 + 77*x1*x4**3/8 + 13*x2**4/12 + 203*x2**3*x3/48 - 367*x2**2*x3**2/32 + 103*x2**2*x4**2/32 - 47*x2*x3**3/24 - 41*x2*x3*x4**2/8 - 71*x3**4/32 + 221*x3**2*x4**2/16 - 229*x4**4/96",
    "57*x1*x2**2*x3/8 - 249*x1*x2*x3**2/16 + 257*x1*x2*x4**2/16 - x1*x3*x4**2 + 10*x2**3*x4/3 + 91*x2*x3**2*x4/4 - 263*x2*x4**3/24 - 5*x3*x4**3/12",
    "5*x1**3*x3/3 - 6*x1**2*x2*x4 - 3*x1**2*x3*x4/2 - 7*x1*x2**2*x3 - 3*x1*x2*x3**2 + 49*x1*x2*x4**2/16 + 17*x1*x3**3/12 - 25*x1*x3*x4**2/8 - x2**3*x4/3 - 3*x2**2*x3*x4/2 - 9*x2*x3**2*x4/4 + 11*x2*x4**3/8 + 7*x3*x4**3/24",
    "49*x1**2*x4**2/32 + 11*x1*x4**3/24 + 3*x2**3*x3 + 3*x2**2*x3**2 - 81*x2**2*x4**2/32 + 3*x2*x3*x4**2/8 - 9*x3**2*x4**2/16 + 3*x4**4/16",
    "-5*x4**4/6",
    "0"
  ],
  "H6": [
    "89*x1**5*x4/320 - 35*x1**4*x2**2/6 + 89*x1**4*x2*x3/64 - 1015*x1**4*x3**2/288 + 1255*x1**4*x4**2/288 - 371*x1**3*x2**2*x4/48 - 745*x1**3*x2*x3*x4/72 - 787*x1**3*x3**2*x4/288 + 1583*x1**3*x4**3/432 - 59*x1**2*x2**4/32 + 14651*x1**2*x2**3*x3/288 - 1121*x1**2*x2**2*x3**2/48 + 25*x1**2*x2**2*x4**2/6 - 35461*x1**2*x2*x3**3/576 + 7973*x1**2*x2*x3*x4**2/192 - 13499*x1**2*x3**4/1152 + 9743*x1**2*x3**2*x4**2/192 - 7415*x1**2*x4**4/1152 + 15107*x1*x2**4*x4/576 + 845*x1*x2**3*x3*x4/72 + 633*x1*x2**2*x3**2*x4/32 - 453*x1*x2**2*x4**3/8 + 9617*x1*x2*x3**3*x4/72 - 2195*x1*x2*x3*x4**3/24 + 129*x1*x3**4*x4/4 - 10201*x1*x3**2*x4**3/288 + 1031*x1*x4**5/720 - 3*x2**6/32 - 18749*x2**5*x3/2880 + 1889*x2**4*x3**2/96 - 1085*x2**4*x4**2/96 + 30503*x2**3*x3**3/1728 - 14207*x2**3*x3*x4**2/576 - 13859*x2**2*x3**4/1152 - 1133*x2**2*x3**2*x4**2/192 + 3971*x2**2*x4**4/384 - 331*x2*x3**5/144 + 341*x2*x3**3*x4**2/36 + 13*x2*x3*x4**4/12 + 1001*x3**4*x4**2/96 - 2905*x3**2*x4**4/144 + 3031*x4**6/1440",
    "701*x1**3*x2**2*x3/48 + 107*x1**3*x2*x4**2/12 - 295*x1**3*x3*x4**2/48 + 8777*x1**2*x2**3*x4/288 + 427*x1**2*x2**2*x3*x4/16 - 3233*x1**2*x2*x4**3/96 + 421*x1**2*x3*x4**3/12 - 11107*x1*x2**4*x3/288 + 1921*x1*x2**3*x3**2/72 - 193*x1*x2**3*x4**2/72 + 3713*x1*x2**2*x3**3/48 - 133*x1*x2**2*x3*x4**2/4 + 511*x1*x2*x3**2*x4**2/16 - 2653*x1*x2*x4**4/144 + 11*x1*x3**3*x4**2/3 - 715*x1*x3*x4**4/48 - 7297*x2**5*x4/1440 + 1055*x2**4*x3*x4/72 + 2149*x2**3*x3**2*x4/96 + 1165*x2**3*x4**3/144 - 1565*x2**2*x3**3*x4/36 - 13*x2**2*x3*x4**3/4 - 113*x2*x3**2*x4**3/32 - 53*x2*x4**5/48 + 2843*x3**3*x4**3/432 - 1103*x3*x4**5/240",
    "-4*x1**5*x3/3 + 10*x1**4*x2*x4/3 - 251*x1**4*x3*x4/192 - 11*x1**3*x2**2*x3/24 - 859*x1**3*x2*x3**2/96 + 389*x1**3*x2*x4**2/48 - 9487*x1**3*x3**3/1728 - 2717*x1**3*x3*x4**2/576 - 19*x1**2*x2**3*x4/24 - 325*x1**2*x2**2*x3*x4/96 - 487*x1**2*x2*x3**2*x4/64 + 11207*x1**2*x2*x4**3/576 + 3023*x1**2*x3**3*x4/144 - 8371*x1**2*x3*x4**3/576 + 97*x1*x2**4*x3/48 + 705*x1*x2**3*x3**2/32 - 967*x1*x2**3*x4**2/48 - 139*x1*x2**2*x3**3/64 + 1437*x1*x2**2*x3*x4**2/64 - 769*x1*x2*x3**4/288 - 161*x1*x2*x3**2*x4**2/32 + 2449*x1*x2*x4**4/576 + 4021*x1*x3**5/960 - 14885*x1*x3**3*x4**2/288 + 2007*x1*x3*x4**4/64 + x2**5*x4/40 - 79*x2**4*x3*x4/192 - 8041*x2**3*x3**2*x4/576 + 11095*x2**3*x4**3/1728 - 27*x2**2*x3**3*x4/16 + 10175*x2**2*x3*x4**3/576 - 361*x2*x3**4*x4/64 + 7325*x2*x3**2*x4**3/288 - 5153*x2*x4**5/576 + 2479*x3**3*x4**3/288 - 13261*x3*x4**5/1440",
    "115*x1**2*x2**3*x3/24 + 6121*x1**2*x2**2*x3**2/192 - 4975*x1**2*x2**2*x4**2/192 + 1681*x1**2*x2*x3*x4**2/96 - 2059*x1**2*x4**4/1152 - 133*x1*x2**3*x3*x4/48 - 149*x1*x2**2*x3**2*x4/32 - 3169*x1*x2**2*x4**3/288 - 19*x1*x2*x3*x4**3/32 + 3799*x1*x4**5/720 - 77*x2**5*x3/48 - 125*x2**4*x3**2/192 + 673*x2**4*x4**2/192 + 2723*x2**3*x3**3/288 - 1193*x2**3*x3*x4**2/48 + 3199*x2**2*x3**4/576 - 755*x2**2*x3**2*x4**2/48 + 25*x2**2*x4**4/128 + 691*x2*x3**3*x4**2/144 - 521*x2*x3*x4**4/72 - 965*x3**2*x4**4/192 - 1931*x4**6/8640",
    "-167*x1**2*x4**4/16 - 511*x1*x4**5/288 - 655*x2**2*x4**4/48 - 5803*x2*x3*x4**4/288 - 1669*x3**2*x4**4/144 + 1133*x4**6/432",
    "0"
  ]
}''')
x=s.symbols('x1:5')
loc=dict(zip(D['variables'],x))
def mat(v):
    a,b,c,d,e,f=v
    return s.Matrix([[a,b,c,d],[b,-a,d,-c],
                     [c,d,e,f],[d,-c,f,-e]])
H2,H4,H6=[mat([s.sympify(v,locals=loc) for v in D[k]])
           for k in ['H2','H4','H6']]
Om=s.Matrix([[0,1,0,0],[-1,0,0,0],
             [0,0,0,1],[0,0,-1,0]])
for H in [H2,H4,H6]:
    assert H.T==H
    assert H*Om+Om*H==s.zeros(4)
g4=(H4+H2**2/2).applyfunc(s.expand)
g6=(H6+(H2*H4+H4*H2)/2+H2**3/6).applyfunc(s.expand)
iv4=(-H4+H2**2/2).applyfunc(s.expand)
def christ(M):
    return {(k,i,j):s.expand((s.diff(M[k,j],x[i])
             +s.diff(M[k,i],x[j])-s.diff(M[i,j],x[k]))/2)
            for k in range(4) for i in range(4) for j in range(4)}
C1,C3,C5=[christ(H) for H in [H2,g4,g6]]
G1=C1
G3={}
G5={}
for k in range(4):
    for i in range(4):
        for j in range(4):
            G3[k,i,j]=s.expand(C3[k,i,j]-sum(
                H2[k,l]*G1[l,i,j] for l in range(4)))
            G5[k,i,j]=s.expand(C5[k,i,j]-sum(
                H2[k,l]*C3[l,i,j] for l in range(4))+sum(
                iv4[k,l]*G1[l,i,j] for l in range(4)))
def deriv(G,i,j):
    return sum(s.diff(G[k,i,j],x[k])-s.diff(G[k,i,k],x[j])
               for k in range(4))
def prod(G,H,i,j):
    return sum(G[k,i,j]*H[l,k,l]-G[l,i,k]*H[k,j,l]
               for k in range(4) for l in range(4))
for i in range(4):
    for j in range(4):
        assert s.expand(deriv(G1,i,j)+int(i==j))==0
        assert s.expand(deriv(G3,i,j)+prod(G1,G1,i,j)+H2[i,j])==0
        assert s.expand(deriv(G5,i,j)+prod(G1,G3,i,j)
                        +prod(G3,G1,i,j)+g4[i,j])==0
a=[]
for k in range(4):
    def cov(i,j):
        return -sum(G1[l,k,i]*Om[l,j]+G1[l,k,j]*Om[i,l]
                    for l in range(4))
    a.append(s.expand((cov(0,2)-cov(1,3))/2))
assert a==[x[0]-x[3],-x[1]-x[2],x[1]/2,x[0]/2]
A=s.Matrix(a).jacobian(x)
assert A.det()==s.Rational(1,4)
print('Exact assertions passed: compatibility, Einstein residual O(r^6),')
print('and isolated positive-index leading torsion derivative.')

\end{lstlisting}
\begin{thebibliography}{9}
\bibitem{Armstrong} J. Armstrong, On four-dimensional almost K\"ahler manifolds, Quart. J. Math. 48 (1997), 405--415. \href{https://doi.org/10.1093/qmath/48.4.405}{doi:10.1093/qmath/48.4.405}.
\bibitem{AD} V. Apostolov and T. Dr\u aghici, The curvature and the integrability of almost-K\"ahler manifolds: a survey, Fields Institute Communications 35 (2003), 25--53. \href{https://arxiv.org/abs/math/0302152}{arXiv:math/0302152}, especially Section 4, Propositions 2--3 and Theorem 3.
\bibitem{LeBrun} C. LeBrun, Einstein metrics, harmonic forms, and symplectic four-manifolds, Ann. Global Anal. Geom. 48 (2015), 75--85. \href{https://www.math.stonybrook.edu/~claude/cake.pdf}{Author's manuscript}, equation (2).
\bibitem{Minimality} C. LeBrun, Einstein metrics, symplectic minimality, and pseudo-holomorphic curves. \href{https://arxiv.org/abs/math/0409290}{arXiv:math/0409290}, Proposition 2.1 and Section 3.
\end{thebibliography}
\end{document}
